% =====================================================================
% Kapitel 1: Einführung und Projektüberblick
% =====================================================================

\chapter{Einführung und Projektüberblick}

\label{ch:einfuehrung}

\section{Das Problem}

Quantitative Finanzforschung lebt von Daten -- aber an genau diesen
Daten scheitern die meisten Hobby- und Kleinprojekte:

\begin{itemize}
  \item \textbf{Datenquellen sind fragmentiert}. Makro-Daten kommen vom
        FED, Aktien-Daten von Yahoo, Anleihe-Daten vom Treasury, FX-Daten
        von der EZB, Rohstoff-Daten von der COMEX. Jede Quelle hat ein
        anderes Format, eine andere API, andere Lizenzbedingungen.
  \item \textbf{Historie ist teuer}. Bloomberg-Terminals kosten
        20.000\,{}/Jahr, Refinitiv vergleichbar. Für
        Hobby-Forscher unbezahlbar.
  \item \textbf{Open Data ist da, aber unhandlich}. FRED, EZB SDMX,
        Yahoo Finance -- alle frei verfügbar, aber jede mit eigener
        API-Linie, eigenen Rate-Limits, eigenem Datenformat.
  \item \textbf{Reproduzierbarkeit ist schwer}. Wenn der Code zur
        Datensammlung nicht versioniert und mitprotokolliert ist,
        weiß man sechs Monate später nicht mehr, welcher Pull
        welche Daten geliefert hat.
  \item \textbf{Backtesting braucht saubere PIT-Daten}. Eine
        Aktienkurs-Zeitreihe vom 1.~Januar 2020 muss zu jedem
        späteren Zeitpunkt exakt so aussehen, wie sie tatsächlich
        an diesem Tag bekannt war -- inklusive aller nachträglichen
        Splits, Dividenden und Restatements.
\end{itemize}

\section{Die Lösung: \code{deephold\_db}}

\code{deephold\_db} ist eine selbst-gehostete, lokal laufende
Finanzmarktdatenbank, die alle genannten Probleme auf eine pragmatische
Art adressiert:

\begin{definition}
\textbf{Retail-Research-Datenbank}
Eine Datenbank, die \emph{für den privaten und Forschungsgebrauch}
konzipiert ist. Sie ist nicht für Hochfrequenzhandel, regulatorische
Compliance oder institutionelle Vermögensverwaltung gedacht. Sie
priorisiert \textbf{einfache Bedienung} und \textbf{kostengünstigen
Betrieb} über \textbf{Latenz} und \textbf{vollständige PIT-Treue}.
\end{definition}

Die Kernidee: \emph{ein Befehl}, \emph{eine Datenbank}, \emph{mehrere
Datenquellen}, \emph{mehrere Frontends}. Konkret bedeutet das:

\begin{enumerate}
  \item \textbf{Ein-Befehl-Setup}. \code{make init} startet Docker,
        installiert Python-Dependencies, legt das Datenbankschema an
        und seedet die ersten Daten. Alles in unter 10~Minuten.
  \item \textbf{Mehrere Datenquellen, einheitliches Schema}. FRED
        (FED-Wirtschaftsdaten), EZB SDMX (EUR-Geldmarkt, EUR-Anleihen,
        EUR-FX) und Yahoo Finance (US-Aktien, ETFs, Indizes) werden über
        Vendor-Adapter in dasselbe Schema gemappt.
  \item \textbf{Mehrere Frontends}.
        \begin{itemize}
          \item \textbf{Jupyter-Notebook} (\code{build\_notebook.py})
                mit Plotly für interaktive Charts.
          \item \textbf{OpenTUI-Explorer} (Bun + React 19) für
                schnelles Durchsuchen der Datenbank im Terminal.
          \item \textbf{SQL direkt} via psql für Ad-hoc-Analysen.
          \item \textbf{Python-API} (SQLAlchemy ORM) für reproduzierbare
                Research-Skripte.
        \end{itemize}
  \item \textbf{Reproduzierbarkeit}. Jede Zeile in der Datenbank trägt
        einen Zeitstempel (\code{ingested\_at}) und eine
        \code{vendor\_id}. Damit lässt sich auch Monate später
        rekonstruieren, woher ein bestimmter Wert stammt.
  \item \textbf{Walk-Forward-Backtest-Framework}. Eine
        vereinfachte Version des AEGIS-Algorithmus aus dem
        quantitativen Finanzpaper von Chakraborty und Singh
        (2026) demonstriert, wie die gesammelten Daten für
        systematische Strategieentwicklung genutzt werden können.
\end{enumerate}

\section{Anwendungsfälle}

\begin{beispiel}
\textbf{Typische Anwendungsfälle}
  \begin{itemize}
    \item \textbf{Portfolio-Analyse}: ``Wie hat sich mein
          60/40-Aktien/Anleihen-Portfolio in den letzten 20~Jahren
          entwickelt?''
    \item \textbf{Makro-Research}: ``Wie korreliert die
          US-Arbeitslosenquote mit S\&P-500-Renditen in den
          letzten 10~Jahren?''
    \item \textbf{Faktor-Analyse}: ``Lohnt sich Value-Investing
          in Small-Caps noch?''
    \item \textbf{Backtest-Vorbereitung}: ``Lade mir 20~Jahre
          S\&P~500- und NASDAQ-Daten, damit ich meinen
          Mean-Reversion-Filter darauf testen kann.''
    \item \textbf{Bildungsprojekte}: ``Ich will meinen Studierenden
          zeigen, wie ein professionelles Datenprodukt
          aussieht.''
  \end{itemize}
\end{beispiel}

\section{Nicht-Ziele}

\begin{warnung}
\rmfamily\itshape
Folgende Funktionen sind \emph{bewusst nicht} Teil dieses Projekts:

\begin{itemize}
  \item Echtzeitdaten oder Intraday-Ticks.
  \item Echtes Point-in-Time mit allen Splits/Dividenden zum
        jeweiligen historischen Zeitpunkt (nur adjusted-close,
        keine nachträglichen Restatements).
  \item Bond-CUSIP-Level-Daten (nur aggregierte Anleihe-ETFs
        oder Makro-Yields).
  \item Multi-Vendor-Reconciliation (Preisvergleiche zwischen
        Datenquellen).
  \item Historische Index-Mitgliedschaften (``Wann wurde AAPL
        in den S\&P~500 aufgenommen?'').
  \item Historische Credit-Ratings.
\end{itemize}
\end{warnung}

\section{Architektur auf einen Blick}

\begin{figure}[ht]
\centering
\begin{tikzpicture}[
  node distance=0.8cm and 1.2cm,
  box/.style={rectangle, draw, thick, rounded corners, minimum width=2.8cm, minimum height=0.8cm, align=center, font=\small},
  layer/.style={rectangle, draw, very thick, rounded corners, fill=#1!10, inner sep=0.4cm},
  arrow/.style={-Latex, thick},
]
  \node[box, fill=blue!20] (python) {Python 3.11\\(SQLAlchemy, Polars)};
  \node[box, fill=green!20, right=of python] (db) {PostgreSQL 16\\(14 Tabellen)};
  \node[box, fill=yellow!20, right=of db] (bun) {Bun + React 19\\(OpenTUI)};

  \node[box, fill=red!20, below=of python] (vendors) {Vendor-Adapter\\(FRED, ECB, Yahoo)};
  \node[box, fill=cyan!20, below=of db] (etl) {ETL-Pipeline\\(Seeder, query\_db.py)};
  \node[box, fill=orange!20, below=of bun] (notebook) {Plotly-Notebook\\(build\_notebook.py)};

  \draw[arrow] (vendors) -- node[above, font=\scriptsize] {insert} (db);
  \draw[arrow] (etl) -- node[above, font=\scriptsize] {read/write} (db);
  \draw[arrow] (python) -- node[above, font=\scriptsize] {SQLAlchemy} (db);
  \draw[arrow] (bun) -- node[above, font=\scriptsize] {node-postgres} (db);
  \draw[arrow] (notebook) -- node[above, font=\scriptsize] {read-only} (db);
\end{tikzpicture}
\caption{High-Level-Architektur von \code{deephold\_db}: Python-Backend
  füllt die DB über Vendor-Adapter, das TUI und das Notebook lesen
  direkt aus der DB. Die einzige Kopplung zwischen den Layern ist
  die Datenbank selbst.}
\label{fig:arch-overview}
\end{figure}

Abbildung~\ref{fig:arch-overview} zeigt die drei Schichten des
Projekts:

\begin{enumerate}
  \item \textbf{Python-Backend} (links): Enthält die Vendor-Adapter
        (FRED, ECB, Yahoo), die ETL-Pipeline (Seeder, \code{query\_db.py})
        und die Analytics-Module (Returns, Metriken, VAM-Backtest).
        Spricht via SQLAlchemy mit der DB.
  \item \textbf{PostgreSQL-Datenbank} (Mitte): Single Source of Truth.
        14~Tabellen -- \code{venues}, \code{vendors}, \code{instruments},
        \code{instrument\_identifiers}, \code{trading\_calendars},
        \code{prices\_raw}, \code{prices\_daily}, \code{corporate\_actions},
        \code{fx\_rates\_daily}, \code{bond\_yields}, \code{macro\_series},
        \code{macro\_observations}, \code{dq\_runs}, \code{dq\_findings}.
  \item \textbf{TypeScript-TUI} (rechts): OpenTUI-Explorer in Bun.
        Liest die DB via \code{node-postgres}, kein Python im Hot-Path.
        Maximale Performance für den Endbenutzer.
\end{enumerate}

\section{Was Sie als Nächstes erwartet}

Das nächste Kapitel~--\ref{ch:grundlagen}--~bringt Sie auf den
nötigen Wissensstand, falls Sie mit einem der verwendeten Konzepte
noch nicht vertraut sind. Falls Sie bereits Erfahrung mit PostgreSQL,
Docker, Python und TypeScript haben, können Sie direkt mit
Kapitel~\ref{ch:infrastruktur} fortfahren.
